Fix a law in the prespecified class \(\mathcal M_{n,d,0}\).  By \({\rm ass:third-derivative-bound}\) with radius zero, \(\partial_x^3f_i(a,x)=0\) for every displayed argument.  Since each response is \(C^3\), the fundamental theorem of calculus gives, for \(x,y\in[0,1]\),
\[
 \partial_x^2f_i(a,x)-\partial_x^2f_i(a,y)
 =\int_y^x\partial_t^3f_i(a,t)\,dt=0.
\]
Thus the second derivative is constant and two further integrations show that \(x\mapsto f_i(a,x)\) is a polynomial of degree at most two.

Substitution of \(L=0\) into \({\rm def:optimal-bandwidth}\) gives
\[
 h_0=\infty,
 \qquad h_\star=\min\{d^{-1/2},\max\{d^{-1},\infty\}\}=d^{-1/2}.
\tag{1}
\]
The logarithmic-degree condition in \({\rm def:model-class}\) implies \(d\to\infty\), so for all sufficiently large indices \(d^{-1}\leq d^{-1/2}\leq1/2\); hence (1) belongs to \(\mathcal H_{n,d}\).  The positive-degree Taylor term in \({\rm def:radii}\) is multiplied by \(L\) and is therefore zero.

Fix \(i\) with \(N_i>0\) and \(a\in\{0,1\}\).  By \({\rm lem:exact-score-moments}\),
\[
 \mathbb E_Z[q_{i,h}\mid A]=0,
 \qquad \mathbb E_Z[q_{i,h}V_i\mid A]=1,
 \qquad \mathbb E_Z[q_{i,h}V_i^2\mid A]=0.
\tag{2}
\]
The exact quadratic expansion and (2) give
\[
 \mathbb E_Z[q_{i,h}f_i(a,1/2+V_i)\mid A]=\partial_xf_i(a,1/2),
\tag{3}
\]
so every quadratic response, and in particular every affine response, has zero positive-degree approximation error.

At \(h_\star=d^{-1/2}\), the analytic stochastic term in \({\rm def:frontier-envelope}\) is
\[
 (n\sqrt d\,h_\star^3)^{-1/2}
 =(n\sqrt d\,d^{-3/2})^{-1/2}=\sqrt{d/n},
\tag{4}
\]
and the graph term is also \(\sqrt{d/n}\).  Direct substitution of \(L=0\) and \(h=h_\star\) into \({\rm def:pc-handle}\) and \({\rm def:radii}\) yields
\[
 \overline b_{h_\star}(A)
 =\frac Bn\sum_{i\in\mathcal V_n}\mathbf1\{N_i=0\}
 +\widehat\Delta_{\mathrm{pc},h_\star},
\tag{5}
\]
\[
 \overline s_{h_\star}(A)^2
 =\widehat V_{\mathrm{pc},h_\star}
 =\frac1{4n^2}\sum_{j\in\mathcal V_n}
 \left(\sum_{i\in\mathcal V_n}\beta_{ij,h_\star}\right)^2.
\tag{6}
\]
The certification in \({\rm thm:explicit-pc-conditional-certificate}\) makes (5)--(6) the complete conditional certificates; no asymptotic order for (6) is used.

All member properties other than the third-derivative radius are unchanged when that radius is reduced from \(L\) to zero, so \(\mathcal M_{n,d,0}\subseteq\mathcal M_{n,d,L}\).  To verify the non-post-hoc qualification when \(L>0\), retain the graph and uniform latent-type law of an admissible member and set deterministically
\[
 g(a,x)=c(x-1/2)^3,
 \qquad c=\frac12\min\{B/3,L/6\}>0.
\tag{7}
\]
This preserves unit sampling, randomization, the graph law, and anonymous interference.  Moreover,
\[
 \|g\|_\infty\leq c/8,\quad
 \|\partial_xg\|_\infty\leq3c/4,\quad
 \|\partial_x^2g\|_\infty\leq3c,\quad
 \|\partial_x^3g\|_\infty=6c,
\]
which are bounded by \(B,B,B,L\), respectively.  Thus this alternative lies in \(\mathcal M_{n,d,L}\) but not in \(\mathcal M_{n,d,0}\).  For \(N_i>0\),
\[
 \mathbb E_Z[q_{i,h}g(a,1/2+V_i)\mid A]
 =c\frac{\mathbb E_Z[V_i^4\mathbf1\{|V_i|\leq h_i(h)\}\mid A]}{\mu_{2,i}(h)}>0.
\tag{8}
\]
The denominator is positive by \({\rm lem:exact-score-moments}\), and the numerator is positive because the truncation contains a nearest nonzero lattice point of positive binomial mass.  Hence a realized affine law inside the positive-radius class cannot justify replacing the prespecified honesty class by the zero-radius class.  Equations (1)--(8) prove the proposition.